%@ pages 68-72
%@ status translated
% Original section 7, "Catégories de Segal et catégories simpliciales", pp. 68--72.

\clearpage
\section{Segal categories and simplicial categories}
\label{segsimplpage}

Most of the examples of Segal $n$-stacks concern the case $n=1$. It is worth
emphasising once again that a Segal $1$-stack is in fact a certain type of
$\infty$-stack, in which the morphisms in degree $\geq 2$ are invertible. In
particular, a Segal $1$-stack may be an $n$-stack for $n>1$ (this is the case
when the simplicial sets of morphisms are $n-1$-truncated).


A {\em simplicial category}
\footnote{We make a slight abuse of notation here: in principle a
``simplicial category'' should rather be a simplicial object in $Cat$, in
particular the objects should form a simplicial set. We impose the additional
condition that the simplicial set of objects be constant, i.e. a discrete
set.}
is a category enriched in simplicial sets (see \cite{Kelly}). It is therefore a
set of objects $Ob(C)$, equipped, for every pair of objects $x,y\in Ob(C)$,
with a simplicial set $Mor_C(x,y)$ said to be that of morphisms from $x$ to
$y$, together with a composition law
$$
Mor_C(x,y)\times Mor_C(y,z)\rightarrow Mor_C(x,z)
$$
which is strictly associative, and with identities $1_x\in (Mor_C(x,x))_0$.

Simplicial categories have for a long time already played a key role in
homotopy theory, starting with Kan and Quillen; see subsequently
\cite{DwyerKan1} \cite{DwyerKan2} \cite{DwyerKan3}, particularly the
introduction of \cite{DwyerKanDiags}, and more recently \cite{DHK}; in the
same volume as \cite{DwyerKanDiags} one immediately notices other papers that
use this notion in an essential way, for example that of Waldhausen.

If $C$ is a simplicial category, one defines the {\em associated Segal
$1$-category}, which is still denoted $C$, as the Segal category defined by
$$
C_0:= Ob(C),
$$
and, for $x_0,\ldots , x_p\in C_0$,
$$
C_{p/}(x_0,\ldots ,x_p):= Mor_C(x_0,x_1)\times Mor _C(x_1,x_2)\times \ldots
\times Mor_C(x_{p-1},x_p).
$$
The morphisms of the simplicial structure are defined (in the obvious way)
using composition where needed. For example, composition itself becomes the
face morphism $02$
$$
C_{1/}(x,y)\times C_{1/}(y,z)=C_{2/}(x,y,z)\rightarrow C_{1/}(x,z).
$$

A functor between two simplicial categories is an equivalence if and only if
the corresponding morphism between Segal $1$-categories is an equivalence. We
shall show, using the theory of Dwyer--Kan localisation which will be treated
below, that every Segal $1$-category is equivalent to a simplicial category.
More precisely, if $A$ is a Segal $1$-category and $A'$ its fibrant
replacement, then there exist a simplicial category $C$ and an equivalence
$C\rightarrow A'$.

For our treatment of morphisms, stacks, etc. we systematically use the point
of view of Segal $1$-categories. Let us point out here that the techniques we
use are also available, or in the process of being developed, for simplicial
categories. Dwyer, Hirschhorn and Kan describe a closed model category for
simplicial categories, in which the cofibrations are the retracts of
(iterated) free extensions \cite{DHK}. On the other hand, Cordier and Porter
have developed a theory of ``homotopy coherent'' morphisms between two
simplicial categories \cite{CordierPorter}. If $C$ and $D$ are two simplicial
categories, they define a simplicial category $Coh(C,D)$ of ``coherent''
functors and natural transformations between $C$ and $D$. The notion of
composition for such functors is a little problematic (because the Segal
$2$-category $1SeCAT$ is not equivalent to a strict simplicial
$2$-category), and this is the subject of a large part of the paper
\cite{CordierPorter}.

There notably remains the task of making the link between these different
points of view. One can isolate two problems:


\begin{problem}
\label{probleme1}
Give a (Quillen) equivalence between the cmc of simplicial categories of
Dwyer--Hirschhorn--Kan \cite{DHK}, and the cmc of Segal $1$-precategories of
Theorem \ref{SeCmf}; % the French prints "1-pr\'ecats" here
\end{problem}

%%%%%%% added in v3 %%%%%%
The problem above is to a large extent
treated in Dwyer--Kan--Smith \cite{DKS}:
they give an equivalence between the homotopy categories; it remains to
verify that their construction gives a Quillen functor.
%%%%%%%%%%

\begin{problem}
\label{probleme2}
Given two simplicial
categories $C$ and $D$, construct an equivalence of
Segal $1$-categories
$$
Coh (C,D) \cong \underline{Hom}(C',D')
$$
where $Coh(C,D)$ is the simplicial category defined by Porter and Cordier
\cite{CordierPorter},
$C'$ and $D'$ are fibrant replacements of the
Segal categories corresponding to $C$ and $D$,
and $\underline{Hom}(C',D')$ is the internal $Hom$ of Segal $1$-precategories
(which will be treated in \S 11 below).
\end{problem}

We take a step towards the first of these problems in Corollary
\ref{catsimpl} below, where we show that every Segal $1$-category is
equivalent to a simplicial $1$-category.

{\em Remark:} The closed model category of Dwyer--Hirschhorn--Kan
cannot be ``internal'' (see \S 11 below) because the cofibrations
are (iterated) free extensions of categories and are therefore not stable
under direct product. In particular, the construction of \cite{nCAT}
(see also \S 11 below) of a strict enriched category for the
objects in question cannot be adapted to
this closed model category.

In fact, there cannot
exist a strict category $SplCAT$ enriched over simplicial
$1$-categories and having the right homotopy type (i.e.
equivalent to $1SeCAT$). For if $SplCAT$ existed, one could perform
the truncation $\tau _{\leq 3}SplCAT$, that is to say a truncation $\leq 2$
for the simplicial categories of morphisms, or else a truncation
$\leq 1$ for the
simplicial sets of $2$-morphisms. Now the truncation $\tau_{\leq 1}$ gives
a strict $1$-category, and is compatible with direct products. This
would entail that $\tau _{\leq 3}SplCAT$ would be a strict $3$-category.
But $1SeCAT$ has non-trivial Whitehead products, which implies
that $\tau _{\leq 3}(1SeCAT)$
cannot be equivalent to a strict $3$-category.


\subnumero{Homotopy products and coproducts}

In a simplicial category one has a notion of {\em homotopy fibre
product} and the dual notion of {\em homotopy coproduct}. In fact, these
notions are available in a Segal $n$-category, cf \cite{limits},
but we restrict ourselves here to the case $n=1$. This is a particular case
of the notion of limit (or of
colimit), see \cite{limits} and \S 14 below.

Let $A$ be a simplicial category, with $X,Y,Z\in A_0$ and
morphisms (i.e. vertices of the simplicial set $Hom _A(\cdot , \cdot )$)
$f: X\rightarrow Y$ and $g: Z\rightarrow Y$. If $U\in A_0$ is an object
equipped with morphisms $p: U\rightarrow X$ and $q: U\rightarrow Z$ and with a
homotopy
$\gamma : fp\sim gq$, we shall say that $(U,p,q,\gamma )$ is a {\em homotopy
fibre product} and we shall write
$$
U= X\times ^h_YZ
$$
if, for every object $V\in A_0$, the morphism $\mu _{V,U}$ induced by $(p,q,
\gamma )$
$$
Hom _A(V,U)\rightarrow Hom _A(V,X)\times ^h_{Hom _A(V,Y)}Hom _A(V,Z)
$$
is a weak equivalence; on the right this is the homotopy fibre product
of simplicial sets (obtained by replacing one of the morphisms
by a Kan fibration), and the morphism is defined up to homotopy
after replacement on the right by a Kan simplicial set. In fact
the reader might imagine that $fp=gq$ and that $\gamma$ is the constant
homotopy---this will be the case in particular in our examples coming from
closed model categories---but if one imposed this condition then
the notion would not be invariant under an equivalence $A\cong A'$ of
simplicial categories. If
$fp$ and $gq$ are equal, the morphism above is well defined with values in the
usual strict (non-homotopy) fibre product.

 If $(V, p',q',\gamma ')$ happens to be another homotopy fibre product
(for the same data)
then the homotopy fibre of $\mu _{V,U}$ over $(p',q',\gamma ')$ is
contractible, just as that of $\mu _{U,V}$
over $(p,q,\gamma)$, which
provides canonical contractible simplicial
sets of morphisms $U\rightarrow V$ and
$V\rightarrow U$, and ensures the essential uniqueness of the homotopy fibre
product $(U,p,q,\gamma )$ if one exists.

We shall say that {\em $A$ admits homotopy fibre products} if there
exists a
homotopy fibre product $(U,p,q,\gamma )= X\times ^h_YZ$ for every
pair of morphisms $X\rightarrow Y\leftarrow Z$ of $A$.

One has a dual notion of {\em homotopy coproduct}: if
$X\stackrel{f}{\leftarrow} Y\stackrel{g}{\rightarrow} Z$ is a
diagram in
$A$, the corresponding homotopy coproduct is a
quadruple $(U, i,j,\delta )$ with $i: X\rightarrow U$,
$j: Z\rightarrow U$ and $\delta$ a homotopy $if\sim jg$. We leave it to the
reader
to give details analogous to what was said above for the product.
We shall denote the homotopy coproduct by
$$
(U,i,j,\delta ) = X\cup _h^YZ,
$$
and we shall say that $A$ {\em admits homotopy coproducts} if there exists
such a homotopy
coproduct for every pair of morphisms
$X\stackrel{f}{\leftarrow} Y\stackrel{g}{\rightarrow} Z$.

These notions are the generalisations to the objects of an arbitrary
simplicial category of the $holim$ and $hocolim$ of Bousfield--Kan
\cite{BousfieldKan} {\em et al.}

These notions can be defined in the same way
for a Segal $1$-category $A$, on condition that one
chooses a splitting $A_{1/}\times
_{A_0}A_{1/}\rightarrow A_{2/}$ (such a splitting exists for example
if $A$ is fibrant)
providing a composition morphism
$A_{1/}\times
_{A_0}A_{1/}\rightarrow A_{1/}$; and one does not need the higher coherence
relations. We leave this development
to the reader, as well as the task of verifying that these
notions of limits and colimits are compatible with the definitions of
\cite{limits}.

In \cite{GabrielZisman}, Gabriel and Zisman introduced a notion of homotopy
fibre product in a strict $2$-category, and they used it
to construct the functor $\Omega$ on the homotopy category $Ho(Top)$
giving rise to the long exact sequence of a fibration. Dually they
introduced a
notion of homotopy coproduct and used it to construct the
functor $\Sigma$
on $Ho(Top)$. In both cases, using the structure of a
$2$-category, they
produce well-defined limits or colimits in the
$1$-category obtained by truncation (in this case the homotopy
category). One can see the notion of homotopy fibre product above
as a generalisation of their notion to the case of
$n$-categories ($1$-groupoidal)
% the French prints "($1$-groupiques)" here
as $n$ tends to infinity. Let us explain in what way the two constructions
are compatible:
if $A$ is a simplicial category then the construction of
\cite{GabrielZisman}
for $Top$ extends to a version for $\tau _{\leq 2}A$, which is a
strict $2$-category; and one can then apply the argument of
\cite{GabrielZisman} to
$\tau _{\leq 2}A$. If $A$ admits homotopy fibre products in the sense
above, then $\tau _{\leq 2}A$ satisfies conditions (A), (B),
(C)
of \cite{GabrielZisman},
and the product they define is the image of our homotopy
fibre product in the truncation $\tau _{\leq 1}A$
(in fact they are mainly interested in
the homotopy fibre over a base point, but for the fibre
product the argument is the same).

The argument of Gabriel--Zisman allows---passing through $\tau _{\leq 2}A$---
one to equip $\tau _{\leq 1}A$
with an additional structure with the
operations $\Omega$ (when $A$ is pointed and admits
homotopy fibre products) and $\Sigma$ (when $A$ is pointed and admits
homotopy coproducts). Alternatively
(but we emphasise all the same that the argument of Gabriel--Zisman
dates from 1964) one can induce directly towards
the truncation $\tau _{\leq 1}A$ these operations, which one can define
on $A$ by
$$
\Omega X := \ast \times ^h_X \ast
$$
and
$$
\Sigma Y:= \ast \cup _h^Y \ast .
$$
A version of the argument of Gabriel--Zisman adapted to the d.g. categories
of complexes is given by Bondal--Kapranov
\cite{BondalKapranov}. We did not become aware of the existence of
the argument of Gabriel--Zisman until M.~Kontsevich pointed out to us
the result of Bondal--Kapranov.
